% --- Configuration de la Langue, de la Police et de l'Encodage ---
\usepackage[french]{babel}
\usepackage{csquotes}
\usepackage{graphicx}
\usepackage{grffile}
% When using babel or polyglossia with biblatex, loading csquotes is recommended to ensure that quoted texts are typeset according to the rules of your main language.

% --- Géométrie (Marges) ---
\usepackage[
    a4paper,
    top=1.5cm,
    bottom=1.75cm,
    left=1.5cm,
    right=1.5cm,
    includehead,
    nomarginpar,
    headheight=0pt
    ]{geometry}

% --- Structure et Table des Matières ---
\usepackage{titlesec} % Pour personnaliser l'apparence des titres de section

\usepackage{tocloft} % Pour personnaliser la table des matières
\setlength{\cftsecnumwidth}{2.5em} % Assure un espace suffisant pour les numéros de section


% --- En-têtes et Pieds de Page ---
\usepackage{fancyhdr}
\pagestyle{fancy}
\fancyhead{} % Vider l'en-tête existant
\fancyhead[L]{\nouppercase{\leftmark}} % Titre de la section actuelle à gauche
\fancyhead[R]{Rapport Année 0} % Titre du rapport à droite
\fancyfoot[C]{\thepage} % Numéro de page centré
\renewcommand{\headrulewidth}{0.4pt} % Ligne sous l'en-tête
\renewcommand{\footrulewidth}{0pt} % Pas de ligne sous le pied de page


% --- Bibliographie (biblatex) ---
\usepackage[
    style=ieee, % Style de citation (Auteur, Année)
    backend=biber,    % Moteur de backend
    sorting=nyvt,     % Tri : Nom, Année, Titre, Volume
    natbib=true,      % Pour utiliser les commandes natbib comme \citep et \citet
]{biblatex}
\addbibresource{source.bib}

\usepackage{amsmath}   % Pour les formules mathématiques
\usepackage{amssymb}
\usepackage[usenames,dvipsnames,svgnames,table]{xcolor}
\usepackage{tikz}
\usepackage{caption}

\usepackage{anyfontsize}
\usepackage{marvosym}

% --- Hyperliens (Doit être le dernier package de la partie structure) ---
\usepackage{hyperref}
\hypersetup{
    colorlinks=true,
    linkcolor=blue,
    filecolor=magenta,
    urlcolor=blue,
    citecolor=green,
}

\baselineskip=18pt  % Interligne de base
\parskip=10pt % Interligne supplémentaire entre paragraphe


% --- Définitions ou redéfinitions de commandes utiles ---
\renewcommand{\thesection}{\Roman{section}}
\renewcommand{\thesubsubsection}{\thesubsection.\alph{subsubsection}}
\newcommand{\todo}[1]{\textbf{\color{red}\uppercase{#1}}}

\titleformat{\section}{\large\bfseries}{\thesection.}{0.5em}{}
\titlespacing*{\section}{0pt}{1.5ex plus 1ex minus .2ex}{1ex plus .2ex}
